% sgd_vs_adamw.tex
% Section analyzing generalization across fine-tuning optimizers

\subsection{Generalization Across Fine-tuning Optimizers}
\label{sec:sgd_vs_adamw}

A critical question for the practical applicability of mergeability prediction is whether conclusions drawn from one training configuration generalize to others. To investigate this, we repeat our full experimental pipeline on models fine-tuned with SGD instead of AdamW, examining whether our L1-regularized linear framework and key metric insights transfer across optimizers.

\paragraph{Experimental Setup.}
We fine-tune ViT-B-16 models on all 20 tasks using SGD with momentum (lr=$10^{-5}$, weight decay=0.1), following the same protocol as our AdamW experiments. We then compute all mergeability metrics and evaluate using our complete suite of prediction methods: L1 Pearson LOTO, standard Pearson LOTO, MSE LOTO, backward elimination, and MLP-based prediction.

\paragraph{Core Findings Hold.}
Table~\ref{tab:sgd_adamw_methods} presents the average validation correlation across all five merge methods for each prediction approach. The fundamental hierarchy of methods remains consistent: Pearson-based objectives substantially outperform MSE optimization, and linear models outperform MLPs.

\begin{table}[h]
\centering
\caption{Average validation Pearson correlation ($r$) across merge methods for different prediction approaches. The ranking of methods is consistent across optimizers.}
\label{tab:sgd_adamw_methods}
\begin{tabular}{lccl}
\toprule
\textbf{Method} & \textbf{SGD} & \textbf{AdamW} & \textbf{Consistent?} \\
\midrule
Pearson (no L1) & \textbf{0.655} & 0.525 & \checkmark \\
L1 Pearson & 0.590 & \textbf{0.544} & \checkmark \\
MLP & 0.450 & 0.493 & \checkmark \\
Backward Elim. & 0.410 & 0.459 & \checkmark \\
MSE & 0.064 & 0.147 & \checkmark \\
\bottomrule
\end{tabular}
\end{table}

Notably, SGD-finetuned models exhibit \emph{higher} predictability with linear methods (0.59--0.66) compared to AdamW (0.52--0.54), suggesting that SGD fine-tuning may produce weight configurations with more linearly separable mergeability characteristics. Interestingly, while L1 regularization improves generalization for AdamW (0.544 vs.\ 0.525), SGD models achieve their best performance \emph{without} regularization (0.655 vs.\ 0.590). This suggests that SGD's implicit regularization during fine-tuning reduces the need for explicit L1 penalties during prediction. The MSE objective continues to fail catastrophically on both optimizers, confirming that correlation-based objectives are essential regardless of the fine-tuning procedure.

\paragraph{Metric Importance Transfers.}
Beyond method-level conclusions, we examine whether the same metrics drive predictions across optimizers. Table~\ref{tab:sgd_adamw_metrics} shows the top-10 most important metrics (ranked by $|\text{coefficient}| \times \text{selection frequency}$) for each optimizer.

\begin{table}[h]
\centering
\caption{Top-10 metrics by importance score for L1 Pearson LOTO. Six metrics appear in the top-10 for both optimizers (marked with $\dagger$).}
\label{tab:sgd_adamw_metrics}
\begin{tabular}{clcl}
\toprule
\textbf{Rank} & \textbf{AdamW} & \textbf{Rank} & \textbf{SGD} \\
\midrule
1 & encoder\_gradient\_l2\_dist$^\dagger$ & 1 & task\_vector\_cosine\_sim \\
2 & input\_gradient\_l2\_dist$^\dagger$ & 2 & layerwise\_effective\_rank \\
3 & task\_vector\_l2\_dist & 3 & layerwise\_eff\_rank\_score \\
4 & right\_subspace\_overlap\_bot & 4 & input\_gradient\_l2\_dist$^\dagger$ \\
5 & right\_subspace\_overlap\_top & 5 & interaction\_overlap\_bot$^\dagger$ \\
6 & activation\_cosine\_sim$^\dagger$ & 6 & task\_vector\_dot\_product \\
7 & singular\_value\_overlap$^\dagger$ & 7 & activation\_dot\_product$^\dagger$ \\
8 & interaction\_overlap\_bot$^\dagger$ & 8 & activation\_cosine\_sim$^\dagger$ \\
9 & activation\_dot\_product$^\dagger$ & 9 & encoder\_gradient\_l2\_dist$^\dagger$ \\
10 & input\_gradient\_dot\_product & 10 & singular\_value\_overlap$^\dagger$ \\
\bottomrule
\end{tabular}
\end{table}

Six of the top-10 metrics (60\%) are shared between optimizers, indicating substantial agreement on which model properties matter for mergeability. The shared metrics span multiple categories:
\begin{itemize}
    \item \textbf{Gradient-based}: \texttt{encoder\_gradient\_l2\_distance}, \texttt{input\_gradient\_l2\_distance}
    \item \textbf{Activation-based}: \texttt{activation\_cosine\_similarity}, \texttt{activation\_dot\_product}
    \item \textbf{Subspace-based}: \texttt{interaction\_matrix\_overlap\_bottom\_k}, \texttt{singular\_value\_overlap}
\end{itemize}

\paragraph{Gradient Metrics Remain Predictive.}
While the relative ranking differs, gradient-based metrics appear in the top-10 for both optimizers. For AdamW, \texttt{encoder\_gradient\_l2\_distance} ranks first with 100\% selection frequency. For SGD, the same metric ranks ninth but maintains 80\% selection frequency. The gradient L2 distance metrics consistently carry negative coefficients across both optimizers, confirming the intuition that models with similar gradient landscapes merge more successfully.

\paragraph{Optimizer-Specific Patterns.}
Some differences emerge between optimizers. SGD models show stronger reliance on effective rank metrics (\texttt{layerwise\_effective\_rank} ranks 2nd), while AdamW models favor subspace overlap metrics (\texttt{right\_subspace\_overlap} ranks 4--5th). These differences likely reflect how each optimizer shapes the loss landscape geometry during fine-tuning. Nevertheless, these optimizer-specific metrics complement rather than contradict the shared predictors.

\paragraph{Conclusion.}
Our linear prediction framework generalizes effectively across fine-tuning optimizers. The core finding---that Pearson correlation objectives vastly outperform MSE, and that simple linear models suffice---holds for both AdamW and SGD. While the optimal regularization strength differs (L1 helps AdamW but not SGD), the underlying linear structure of the prediction problem remains consistent. Moreover, 60\% of the top predictive metrics are shared, with gradient-based and activation-based metrics appearing consistently across both configurations. These results suggest that mergeability is governed by fundamental geometric properties of fine-tuned models that transcend specific optimization choices.
